\documentclass[slidestop]{beamer}

\usepackage[english]{babel}

\title{Dutch Variant Database: design and implementation}
\providecommand{\myConference}{Lab-J work discussion}
\providecommand{\myDate}{Wednesday, 20 June 2012}
\author{Martijn Vermaat}
\providecommand{\myGroup}{}
\providecommand{\myDepartment}{Department of Human Genetics}
\providecommand{\myCenter}{Center for Human and Clinical Genetics}
\providecommand{\lastCenterLogo}{
%  \raisebox{-0.1cm}{
%    \includegraphics[height = 1cm]{lgtc_logo}
%  }
}
\providecommand{\lastRightLogo}{
  \includegraphics[height = 0.7cm]{nbic_logo}
  %\includegraphics[height = 0.8cm]{nwo_logo_en}
}

\usetheme{lumc}

\begin{document}

% This disables the \pause command, handy in the editing phase.
%\renewcommand{\pause}{}

% Make the title page.
\bodytemplate

% First page of the presentation.
\section{Introduction}
\begin{frame}
 \frametitle{Title of this frame.}
 First line.

 \pause

 Second line.
\end{frame}

\section{Topic2}
\begin{frame}
  \begin{itemize}
    \item Item 1.
    \item Item 2.
    \pause
    \item Item 3.
    \pause
    \begin{itemize}
      \item Subitem 1.
    \end{itemize}
  \end{itemize}
\end{frame}

\begin{frame}
  The second slide in the same section.
\end{frame}

\begin{frame}[fragile]
  And an example of displaying code, mind the [fragile] option.

  \begin{lstlisting}[caption = {Example input}]
    print "Hello"
  \end{lstlisting}
\end{frame}

\section{Questions?}
\lastpagetemplate
\begin{frame}
  \begin{center}
    Acknowledgements:
    \bigskip
    \bigskip

  \end{center}
\end{frame}

\end{document}
